


The Security Portal gives security staff at the campus gate a dedicated interface for verifying approved passes, logging student entry and exit, and tracking who is currently outside campus.

<Info>
  Log in with the `security_admin` username. The Security Portal is
  role-restricted — other admin roles cannot access the gate control functions.
</Info>

\subsection{Dashboard overview}

When you log in, the Security Portal dashboard shows four real-time metrics:

- **Students outside** — total number of students currently not in campus
- **Today's check-ins** — students who have returned to campus today
- **Today's check-outs** — students who have left campus today
- **Sync status** — whether the terminal is connected to the central server

These counters update automatically. The portal polls the summary endpoint every five minutes when the tab is active, and every 30 seconds as a fallback.

\subsubsection{Security summary endpoint}

**Endpoint:** `GET /requests/security/summary`

\begin{lstlisting}[language=bash]
GET https://api.uniz.rguktong.in/api/v1/requests/security/summary
Authorization: Bearer <TOKEN>

\end{lstlisting}

\begin{lstlisting}[language=json]
{
  "success": true,
  "summary": {
    "totalOutside": 14,
    "todayCheckins": 8,
    "todayCheckouts": 22
  }
}

\end{lstlisting}

\subsection{Viewing approved passes}

\subsubsection{Approved outpasses}

Fetch all outpass records to verify a student's approved long-duration leave.

**Endpoint:** `GET /requests/outpass/all`

\begin{lstlisting}[language=bash]
GET https://api.uniz.rguktong.in/api/v1/requests/outpass/all
Authorization: Bearer <TOKEN>

\end{lstlisting}

\begin{lstlisting}[language=json]
{
  "success": true,
  "outpasses": [
    {
      "id": "req-123",
      "studentId": "O210008",
      "studentGender": "M",
      "reason": "Home visit",
      "isApproved": true,
      "fromDay": "2026-02-12T09:00:00Z",
      "toDay": "2026-02-15T18:00:00Z",
      "currentLevel": "approved"
    }
  ]
}

\end{lstlisting}

Only show passes where `isApproved: true` to the gate officer. Pending or rejected passes are not valid for exit.

\subsubsection{Approved outings}

Fetch approved short-duration outing requests.

**Endpoint:** `GET /requests/outing/all`

\begin{lstlisting}[language=bash]
GET https://api.uniz.rguktong.in/api/v1/requests/outing/all
Authorization: Bearer <TOKEN>

\end{lstlisting}

\subsection{Gate check-out process}

Use this procedure when a student presents their approved pass at the gate to leave campus.


\section{Instruction Step}

  
\section{Instruction Step}

    Enter the student's university ID in the scan field on the Security Portal dashboard. You can type the ID manually or scan the student's ID card.

    \begin{lstlisting}[language=bash]
    POST https://api.uniz.rguktong.in/api/v1/profile/student/search
    Authorization: Bearer <TOKEN>
    Content-Type: application/json
    
\end{lstlisting}

    \begin{lstlisting}[language=json]
    {
      "username": "O210008",
      "limit": 5
    }
    
\end{lstlisting}

    The response returns matching student records with name, branch, and year.

  
  
\section{Instruction Step}

    Confirm that the student has an approved outpass or outing by checking `isApproved: true` on their request record. Cross-check the `fromDay`/`toDay` (outpass) or `fromTime`/`toTime` (outing) against the current date and time.
  
  
\section{Instruction Step}

    Call the check-out endpoint with the student's username as the request ID. This marks the student as outside campus in the system.

    **Endpoint:** `POST /requests/:id/checkout`

    \begin{lstlisting}[language=bash]
    POST https://api.uniz.rguktong.in/api/v1/requests/O210008/checkout
    Authorization: Bearer <TOKEN>
    
\end{lstlisting}

    \begin{lstlisting}[language=json]
    {
      "success": true,
      "msg": "Checked-Out"
    }
    
\end{lstlisting}

    The student's `is_in_campus` status changes to `false` in the system.

  
</Steps>

\subsection{Gate check-in process}

When a student returns to campus, log their entry to update their presence status.


\section{Instruction Step}

  
\section{Instruction Step}

    Enter the student's university ID in the scan field.
  
  
\section{Instruction Step}

    Call the check-in endpoint to mark the student as back on campus.

    **Endpoint:** `POST /requests/:id/checkin`

    \begin{lstlisting}[language=bash]
    POST https://api.uniz.rguktong.in/api/v1/requests/O210008/checkin
    Authorization: Bearer <TOKEN>
    
\end{lstlisting}

    \begin{lstlisting}[language=json]
    {
      "success": true,
      "msg": "Checked-In"
    }
    
\end{lstlisting}

    The student's `is_in_campus` status changes to `true`.

  
</Steps>

\subsection{Manually updating student presence status}

If you need to correct a student's campus status outside of the normal check-in/check-out flow — for example, if a system error left the status incorrect — use the presence status endpoint directly.

**Endpoint:** `PUT /profile/student/status`

\begin{lstlisting}[language=bash]
PUT https://api.uniz.rguktong.in/api/v1/profile/student/status
Authorization: Bearer <TOKEN>
Content-Type: application/json

\end{lstlisting}

Set `isPresent` to `false` when the student is leaving, and `true` when they have returned.

\begin{lstlisting}[language=json]
{
  "username": "O210008",
  "isPresent": false
}

\end{lstlisting}

<Warning>
  Manually updating presence status bypasses the outpass verification check.
  Only use this endpoint to correct data errors, not as a substitute for the
  standard check-in/check-out flow.
</Warning>

\subsection{Viewing all students currently outside campus}

Generate a list of all students who have not yet returned to campus. This is useful for end-of-day reconciliation and for identifying overdue returns.

**Endpoint:** `POST /profile/student/search`

\begin{lstlisting}[language=bash]
POST https://api.uniz.rguktong.in/api/v1/profile/student/search
Authorization: Bearer <TOKEN>
Content-Type: application/json

\end{lstlisting}

\begin{lstlisting}[language=json]
{
  "isPresentInCampus": false,
  "limit": 50
}

\end{lstlisting}

**Response:**

\begin{lstlisting}[language=json]
{
  "success": true,
  "students": [
    {
      "username": "O210008",
      "name": "DESU SABER",
      "branch": "CSE",
      "is_in_campus": false
    }
  ],
  "pagination": {
    "page": 1,
    "totalPages": 1,
    "total": 14
  }
}

\end{lstlisting}

Increase the `limit` value if you need to fetch more than 50 records in a single request.

\subsection{Activity logs}

The Security Portal includes an **Activity Logs** tab that shows recent entry and exit events recorded at the gate. Use it to review a history of check-in and check-out actions for audit or dispute resolution.


\section{Deep Technical Analysis}
The underlying system architecture for this module is built upon the \textbf{UniZ Microservices Framework}. 
This design ensures that each functional unit (Auth, Profile, Academics) remains isolated, preventing cascading failures across the university ecosystem.

\subsection{Operational Hardening}
Deployments are managed via K3s (Kubernetes) and containerized using high-performance Alpine Linux Docker images. 
Resource management is critical; for instance, the documentation service itself is provisioned with 2GB of RAM to ensure the responsive rendering of complex documentation graphs and state-management components.

\subsection{Enterprise Security Topology}
Every request entering the UniZ gateway is subjected to an aggressive JWT validation layer. 
Permissions are granular, mapping directly onto the RBAC (Role-Based Access Control) matrix defined in the core system specifications.

\subsection{Data Governance}
Prisma-driven data access ensures type safety and predictable query performance. 
We utilize PostgreSQL connection pooling to handle concurrent requests from the student portal and administrative dashboards simultaneously.
